\section{The six programmes after the stipulated RH theorem}
The following are consequences through existing criteria, not assertions that those criteria supplied Alp\"oge's proof. A structural construction ``does not follow'' means that the cited implication gives no such construction; no formal logical independence claim is intended. General-Hecke extensions below remain conditional on Masullo's stated GRH claim.

\subsection{PRG-WEIL: the quadratic form, with its sign specified}
Use $\widetilde g(s)=\int_0^\infty g(x)x^{s-1}dx$ and $g^\sharp(x)=x^{-1}\overline{g(x^{-1})}$. For multiplicative convolution,
\[
 \widetilde{g*g^\sharp}(s)=\widetilde g(s)\overline{\widetilde g(1-\bar s)}.
\]
The explicit formula in the attached Connes--Consani normalization reads
\[
 \widetilde f(0)-\sum_\rho\widetilde f(\rho)+\widetilde f(1)=\sum_v\mathcal W_v(f).
\]
Therefore, for smooth compactly supported $g$ with $\widetilde g(0)=\widetilde g(1)=0$, the Alp\"oge consequence is
\[
 Q(g)=\sum_\rho|\widetilde g(\rho)|^2\ge0,
 \qquad \sum_v\mathcal W_v(g*g^\sharp)=-Q(g)\le0.
\]
This follows by the same Mellin calculation as \fid{MAS-04}; rapid decay and the zero count justify the sum. The converse is Weil's criterion, in the precise version \cite[Appendix C, Proposition C.1, label \texttt{mainprop}]{cc}. It is positivity of a global quadratic form, not positivity of each local term or of a scaling generator. The analogous character-sector statement follows from stipulated Dirichlet GRH with the corresponding completed local factors.

The source's local archimedean result, \fid{PRG-CC-LOCAL}, is already unconditional: \cite[Introduction, Theorem 1, label \texttt{mainthmintro}]{cc} proves
$W_\infty(g*g^*)\ge\Tr(\vartheta(g)\mathbf S\vartheta(g)^*)\ge0$
for its specified support $[2^{-1/2},2^{1/2}]$ and transform vanishings at $i/2$ and $0$. Here $\mathbf S$ is the Sonin-space projection and the notation is the paper's centered normalization. That local compression is real operator structure, but is not an identity with $B$. A geometric realization of the global form by arithmetic correspondences or intersection theory remains an additional target. The curve comparison has exact home in SGA 5, III B, 1.3.1; no EGA construction of the number-field realization is obtained by the inequality alone.

\subsection{PRG-LI / B-03: a positive sequence and an actual paired trace}
Li's criterion is \cite[\S1, Theorem 1]{li}. Its coefficients are
\[
 \lambda_n=\frac1{(n-1)!}\left.\frac{d^n}{ds^n}
 \bigl(s^{n-1}\log\xi(s)\bigr)\right|_{s=1}
 =\sum_\rho\left[1-\left(1-\frac1\rho\right)^n\right],\qquad n\ge1,
\]
where the latter is the symmetric zero sum. All $\lambda_n$ are positive under the stipulated RH. The following checked special form connects this criterion to the designated $B$ and refines SA23 \S4.1(4).

\begin{proposition}
Unconditionally, put $C=B(1+B/4)^{-1}$ and $P_n(x)=2-2T_n(1-x/2)$, where $T_n$ is the Chebyshev polynomial. Then
\begin{equation}\label{eq:litrace}
 \lambda_n=\Tr P_n(C)=\Tr\bigl[2I-2T_n(I-C/2)\bigr]\quad(n\ge1).
\end{equation}
Under RH, $P_n(C)\ge0$ and $\lambda_n>0$.
\end{proposition}
\begin{proof}
The operator $1+B/4$ is invertible: a diagonal zero would require $\gamma_j^2=-1/4$, corresponding to $\rho=0$ or $1$, neither a zero of $\xi$; the diagonal entries tend to one, so their inverses are uniformly bounded. Thus $C$ is trace class, with entries
$x_j=(\gamma_j^2+1/4)^{-1}=[\rho_j(1-\rho_j)]^{-1}$.
For $a=1-1/\rho$, the reflected zero $1-\rho$ gives $a^{-1}$ and
$a+a^{-1}=2-x$. The Chebyshev recurrence gives
$a^n+a^{-n}=2T_n(1-x/2)$. The two Li terms therefore sum to $P_n(x)$.

Since $P_n(0)=0$, $P_n(C)$ is trace class, and the pair sum converges absolutely because $P_n(x)=O_n(x)$ near zero. Pairing preserves the symmetric height cutoff. This proves \eqref{eq:litrace}; the displayed bracket is one polynomial with canceled constant term, not a subtraction of two separately defined infinite traces. Under RH each $x_j$ lies in $(0,4)$, so $T_n(1-x_j/2)\le1$ and $P_n(C)\ge0$. There are arbitrarily small positive $x_j$, and $P_n'(0)=n^2>0$, giving strict positivity of the trace.
\end{proof}
No trace-class operator can have the unpaired spectrum $\{1/\rho\}$, by reciprocal-zero divergence. That narrow intake statement remains correct. A broader assertion that Li coefficients have no trace-class expressions is refuted by \eqref{eq:litrace}. Their positivity supplies no identification of the finite automorphic determinants with this $B$.

\subsection{PRG-NB: the precise closure assertion is now true}
In $L^2((0,\infty),dx)$ put $\rho_a(x)=\{1/(ax)\}$ and $\chi=1_{(0,1]}$. B\'aez-Duarte's strengthened Nyman--Beurling criterion, \cite[Theorem 1.1]{baez}, gives the Alp\"oge consequence
\[
 \chi\in\overline{\operatorname{span}\{\rho_a:a\in\mathbb N\}}^{L^2}.
\]
This is membership of the indicated function, not density of that span in all of $L^2(0,\infty)$. There is no remaining RH gate in this equivalent statement. Prescribed approximants, sharper rates or an arithmetic geometric explanation require their own results; the cited paper already provides constructive approximation information. Cote 1 supplies no identification of these dilations with its determinant operator, and the closure statement itself is functional analysis rather than an EGA/FGA assertion.

\subsection{PRG-DB-FAMILY: de Branges, with the entire shift family}
For every $\varepsilon>0$ set
\[
 E_\varepsilon(z)=\xi(1/2+\varepsilon-iz)=\Xi(z+i\varepsilon),\qquad
 E_\varepsilon^\#(z)=\overline{E_\varepsilon(\bar z)}=\Xi(z-i\varepsilon).
\]
The Alp\"oge consequence is the Hermite--Biehler property
\begin{equation}\label{eq:HB}
 |E_\varepsilon(z)|>|E_\varepsilon^\#(z)|\quad(\Im z>0)
\end{equation}
for \emph{every} $\varepsilon>0$. A numbered family criterion in this convention is \cite[Proposition 1.2; (1.3), (1.7), (1.11)]{suzuki}.
\begin{proof}[Proof of this shifted-$\Xi$ form]
Under RH, $\Xi(z)=\Xi(0)\prod_{\gamma>0}(1-z^2/\gamma^2)$. If $z=x+iy$, $y>0$, each real zero satisfies
$|z+i\varepsilon-\gamma|^2-|z-i\varepsilon-\gamma|^2=4y\varepsilon>0$, and likewise for $-\gamma$. Each paired factor has strictly larger modulus at the upper argument. If the lower-argument product vanishes, the inequality is immediate since the upper argument cannot be a real zero. Otherwise the ratios of partial paired products converge, and fixing one factor bounds every later ratio below a number strictly greater than one. This proves \eqref{eq:HB}. There are no real zeros of $E_\varepsilon$.

Conversely, if all shifts satisfy \eqref{eq:HB} and $\Xi(z_0)=0$ with $\Im z_0>0$, choose $0<\varepsilon<\Im z_0$. Then $E_\varepsilon(z_0-i\varepsilon)=0$ in the upper half-plane, contradicting \eqref{eq:HB}. Conjugation symmetry excludes lower-half-plane zeros as well.
\end{proof}
Each shift consequently defines a de Branges space and a positive reproducing kernel
\[
 K_\varepsilon(w,z)=
 \frac{E_\varepsilon(z)\overline{E_\varepsilon(w)}-
 E_\varepsilon^\#(z)\overline{E_\varepsilon^\#(w)}}{2\pi i(\bar w-z)}.
\]
Positivity here is positivity of all finite kernel Gram matrices. Abstract canonical-system existence follows from de Branges reconstruction; the separate target is a prescribed arithmetic Hamiltonian or kernel, not bare existence. Suzuki's particular arithmetic construction has an explicit parameter range (Theorem 2.3 and \S5). A single large shift is insufficient for RH.

Older sufficient shift-inner-product positivity proposals must be distinguished from this family criterion: Conrey--Li \cite[\S3.1, (3.1)--(3.4)]{conreyli} refutes the proposed $\xi$ choices. Those hypotheses are false, not surviving open consequences of RH. Neither the correct kernel criterion nor those false conditions construct an arithmetic motive or crystal.

\subsection{PRG-CONNES: the RH-equivalent cutoff formula also becomes true}
Connes's local and semilocal trace formulas are established results; the separate global cutoff asymptotic is
\begin{equation}\label{eq:connes}
 \Tr(Q_\Lambda U(h))=2h(1)\log'\Lambda+
 \sum_v\int_{k_v^\times}'\frac{h(u^{-1})}{|1-u|}\,d^\times u+o(1),
\end{equation}
with his cutoffs, measures and test functions \cite[\S VIII, (16), Theorem 5]{connes}. Theorem 5 is first printed in positive characteristic. Its number-field analogue uses the prolate cutoffs in \S VIII, (31)--(37), and the paragraph immediately after (37), which explicitly supplies that analogue. With this qualification, the $\zeta$ sector of \eqref{eq:connes} is true under stipulated RH, and the Dirichlet sectors under stipulated Dirichlet GRH. General-number-field Hecke sectors remain conditional on Masullo here.

Thus the truth of this RH-equivalent asymptotic is not a remaining RH problem. A new direct geometric proof, an identification with \fid{B-01}, and an ordinary trace of the \emph{uncut} scaling action are separate matters. Its relevant positivity is the Weil quadratic distribution, not positivity of the generator. The Lefschetz analogy is with SGA 5's fixed-point traces; the actual adelic noncommutative space is not thereby an \`etale cohomology theory of $\Spec\Z$.

\subsection{PRG-DEN: the cohomology is additional structure}
Deninger's proposed determinant is zeta-regularized:
\[
 \widehat\zeta_K(s)=\prod_{i=0}^2
 \det_\infty\!\left(\frac{s-\Theta}{2\pi}\mid
 H^i_{\mathrm{dyn}}(\overline{\Spec\mathcal O_K},\R)\right)^{(-1)^{i+1}}.
\]
The exact proposal and regularization are \cite[\S3, (2)--(3)]{deninger}; the distributional trace, compact-support version, duality and $K$-theory comparison are (5)--(9). The proposed positive cup/star pairing gives $\Theta+\Theta^*=I$ on degree one, in the common-domain form sense; see pp.~8--9 and the general star-flow relation (11). These are more than zero location.

\paragraph{PRG-DEN-MODEL (what RH alone constructs).}
Using the full zero multiset under RH, on $\ell^2(Z)$ define
\[
 Ae_\rho=\gamma_\rho e_\rho,\quad
 D(A)=\{v:\sum_\rho\gamma_\rho^2|v_\rho|^2<\infty\},\quad
 \Theta=\tfrac12I+iA.
\]
This real diagonal $A$ is selfadjoint: the diagonal inverses of $A\pm i$ are bounded and map onto its domain, or directly its adjoint has the same square-summability domain. Hence $\Theta+\Theta^*=I$ on $D(A)$ and its eigenvalues are the stipulated zeros. This proof establishes a spectral model, not the proposed cohomology. It gives no functorial site, cup product, cycles, duality, or $K$-theory comparison, and no new regularized-determinant theorem. Those remain distinct construction targets. The analogous degree-one Dirichlet spectral statement follows from Alp\"oge; extending to all arithmetic schemes does not follow from RH or Dirichlet GRH. SGA 5 supplies a finite-field model, not an EGA/FGA existence theorem for this number-field object.

\section{PRG-FF: the finite-field dictionary, checked on a curve}
Let $C/\F_q$ be smooth, proper and geometrically connected of genus $g$, with $\ell\nmid q$. Write $V=H^1_{\mathrm{et}}(C_{\overline{\F}_q},\Q_\ell)$ and choose the Frobenius convention with eigenvalues $1$ on $H^0$ and $q$ on $H^2$. Then
\begin{align}
 \#C(\F_{q^n})&=1+q^n-\Tr(F^n\mid V),\label{eq:curve_trace}\\
 Z(C,T)&=\frac{\det(1-TF\mid V)}{(1-T)(1-qT)},\label{eq:curve_det}\\
 -\log\det(1-TF\mid V)&=\sum_{n\ge1}\frac{\Tr(F^n\mid V)}{n}T^n.\label{eq:curve_log}
\end{align}
Exact source locators are SGA 5, III B, Corollary 1.3.1, printed p.~143, and Deligne, Weil I, (1.5.1), (1.5.3), (1.5.4), pp.~275--276 \cite{sga5,weil1}. Equation \eqref{eq:curve_log} is finite linear algebra, so no infinite-dimensional trace theorem is needed. The exterior-power coefficients match cote 1 and SGA 6, V, \S2 and (7.1.1).

Purity, \cite[Theorem (1.6)]{weil1}, says $|\alpha_j|=q^{1/2}$ under each complex embedding. Thus in a complex diagonal realization of the semisimplified eigenvalue multiset, $U=q^{-1/2}F$ is unitary and
\[
 P_1(T)=\det(1-\sqrt q\,T U).
\]
This numerical scaling is not a canonical half Tate twist in the ordinary $\ell$-adic category. Nor is an arbitrary identification of an $\ell$-adic vector space with a complex Hilbert space being asserted. The correct analogue is fixed modulus of eigenvalues, not positive selfadjointness of Frobenius.

\begin{proposition}
For $E/\F_3:y^2=x^3-x$, the curve is smooth, $P_1(T)=1+3T^2$, and $U$ has eigenvalues $i,-i$. No nonzero scalar multiple of its Frobenius can be positive.
\end{proposition}
\begin{proof}
The projective cubic is $Y^2Z-X^3+XZ^2=0$. Its $X$ derivative is $Z^2$ in characteristic three, nonzero on the affine chart; at the sole point at infinity $[0:1:0]$, its $Z$ derivative is one. It is therefore a smooth genus-one curve. For each $x=0,1,2$, $x^3-x=0$, giving three affine points plus infinity. Hence $\#E(\F_3)=4$, the Frobenius trace is $3+1-4=0$, and its determinant on $H^1$ is $3$, since $\Lambda^2H^1\simeq\Q_\ell(-1)$. Thus $P_1(T)=1+3T^2$. The two eigenvalues are $\pm i\sqrt3$. Any nonzero scalar multiple still has opposite nonzero eigenvalues, which cannot both be nonnegative real. Similarity cannot change this obstruction.
\end{proof}
In particular $\det(1-wU)=1+w^2$ and $\Tr U^2=-2$, despite purity. The cup pairing on $H^1$ is alternating with target $\Q_\ell(-1)$; Frobenius is a $q$-similitude. Even that relation alone does not imply purity: $D=\diag(1,q)$ obeys $D^{\mathsf T}JD=qJ$ for $J=\left(\begin{smallmatrix}0&1\\-1&0\end{smallmatrix}\right)$, but its eigenvalue moduli are $1$ and $q$. Direct multiplication proves the claimed identity. Polarization positivity, weight-one purity and positivity of an operator are different statements.

\begin{center}\small
\begin{tabular}{p{0.23\textwidth}p{0.33\textwidth}p{0.33\textwidth}}
\toprule
Feature & $\Xi$ construction & Curve over $\F_q$\\\midrule
Object & Zero-based normal $B$ on $\ell^2$ & Frobenius on finite $H^1_{\mathrm{et}}$\\
Determinant & $\det(1-wB)=G(w)$ & $P_1(T)=\det(1-TF)$\\
Power traces & $\sum_j\gamma_j^{-2n}$ & $1+q^n-\#C(\F_{q^n})$\\
RH statement & $B\ge0$ for this diagonal realization & $|\alpha_j|=\sqrt q$; normalized spectral model unitary\\
Missing geometric datum & Arithmetic realization independent of zero input & Actual sheaf cohomology already supplies $V$, duality and trace\\\bottomrule
\end{tabular}\end{center}
